%% Section 3 -- The CARTA Architectural Framework
%% Ported from docs/DRAFT_PAPER.md Section 3. Math ported verbatim from the
%% draft, which in turn follows docs/PAPER_BLUEPRINT.md 2.2-2.4 (locked).
%%
%% FOUR LOCKED CORRECTIONS -- do not regress any of these:
%%   1. F_1 / F_2 overlap is a feature (delta, epsilon >= 0), not disjointness.
%%   2. ~_dedup is a SIMILARITY relation feeding RECONCILIATION CLUSTERS. No
%%      "equivalence class" or quotient notation anywhere.
%%   3. Idempotence is a softened specification-level target, not a proven
%%      equation S(U,U) = (U, Omega).
%%   4. Normalization has THREE outcomes, not two.

\section{The CARTA Architectural Framework}
\label{sec:carta}

CARTA specifies how to get from a set of fragmented exports to one reconstructed
history together with an account of how that history was produced. It is a
specification and not a system: it fixes what each stage receives, what each
stage emits, and what the result must be able to report about its own
construction. Section~\ref{sec:implementation} describes the planned
implementation, and nothing in this section depends on that implementation
existing.

\subsection{The reconstruction algebra}
\label{sec:carta:algebra}

Write the raw fragment exports as $\mathcal{F} = \{F_1, F_2, \dots, F_n\}$. CARTA
composes three stages (Figure~\ref{fig:carta_pipeline}):

\begin{equation}
  \mathcal{F}
  \xrightarrow{\;\mathcal{I}\;} \mathcal{R}
  \xrightarrow{\;\mathcal{N}\;} (\mathcal{R}', \mathcal{Q})
  \xrightarrow{\;\mathcal{S}\;} (\mathcal{U}, \Omega)
  \label{eq:algebra}
\end{equation}

%% Figure 2 -- CARTA pipeline. The stage signatures are those of Sections 3.2
%% to 3.5, and the caption's contract is the \carta macro from main.tex, so
%% the diagram cannot drift from the text. Two structural facts the prose
%% makes explicit are drawn rather than described: Ingestion and Normalization
%% run per fragment with no cross-fragment reconciliation (3.2), and Synthesis
%% is the only stage that sees fragments together (3.4). The Q path is drawn
%% bypassing Synthesis into Omega, which is what 3.5's quarantine carryover
%% says: quarantined records are carried into the report, not reconciled.
\begin{figure}[htbp]
  \centering
  \begin{tikzpicture}[
      >={Stealth[length=2mm]},
      stage/.style={draw, thick, rounded corners=2pt, minimum height=9mm,
                    minimum width=17mm, align=center, font=\small,
                    inner xsep=1pt},
      data/.style={align=center, font=\small},
      note/.style={font=\scriptsize\itshape, align=center, text=black!65},
      node distance=9mm and 6mm,
    ]

    %% ---- per-fragment lane: Ingestion then Normalization ------------------
    \node[data] (F) {$\mathcal{F}$\\[-1pt]
      {\scriptsize$\{F_1,\dots,F_n\}$}\\[-2pt]
      {\scriptsize$n \ge 2$}};
    \node[stage, right=of F]        (I) {Ingestion\\$\mathcal{I}$};
    \node[data,  right=of I]        (R) {$\mathcal{R}$};
    \node[stage, right=of R]        (N) {Normalization\\$\mathcal{N}$};
    \node[data,  right=of N]        (Rp) {$\mathcal{R}'$};
    \node[stage, right=of Rp]       (S) {Synthesis\\$\mathcal{S}$};
    \node[data,  right=of S]        (U) {$\mathcal{U}$};

    \draw[->] (F) -- (I);
    \draw[->] (I) -- (R);
    \draw[->] (R) -- (N);
    \draw[->] (N) -- (Rp);
    \draw[->] (Rp) -- (S);
    \draw[->] (S) -- (U);

    %% ---- the quarantine path ---------------------------------------------
    \node[data, below=10mm of Rp] (Q) {$\mathcal{Q}$\\[-2pt]
      {\scriptsize$\bigcup_k Q_k$}};
    \draw[->] (N.south) |- (Q.west);

    \node[stage, below=10mm of U, minimum width=22mm] (Om)
      {Report\\$\Omega$};
    \draw[->] (Q) -- (Om)
      node[midway, above, note] {carried, not\\reconciled};
    \draw[->] (S.south) -- ++(0,-5mm) -| (Om.north);

    %% ---- scope annotations, both stated in the prose ---------------------
    \begin{scope}[on background layer]
      \node[draw, dashed, black!45, rounded corners=3pt,
            fit=(I)(N), inner sep=3.5mm] (perfrag) {};
      \node[draw, dashed, black!45, rounded corners=3pt,
            fit=(S), inner sep=3.5mm] (joint) {};
    \end{scope}
    \node[note, above=0.5mm of perfrag, text width=48mm]
      {per fragment; no cross-fragment reconciliation};
    \node[note, above=0.5mm of joint] {fragments\\meet here};

  \end{tikzpicture}
  \caption{The CARTA reconstruction contract \carta. Ingestion and
  Normalization operate on one fragment at a time and perform no
  reconciliation across fragments (Sections~\ref{sec:carta:ingestion}
  and~\ref{sec:carta:normalization}); Synthesis is the only stage that sees
  fragments together (Section~\ref{sec:carta:synthesis}). Records that fail to
  parse, and records that parse under more than one reading, both leave
  Normalization as $\mathcal{Q}$ and are carried into $\Omega$ with their
  reason codes rather than being resolved or discarded
  (Section~\ref{sec:carta:report}).}
  \label{fig:carta_pipeline}
\end{figure}

The quarantine set $\mathcal{Q}$ appears in the top-level algebra by design.
Records that cannot be interpreted safely are an output of the pipeline rather
than an internal concern of one stage, and writing them into the signature is
what makes it impossible for a conforming implementation to discard them quietly.

TCF as defined in Section~\ref{sec:tcf} is a two-way split: one primary device
pauses while one secondary continues, and Definition~1 formalises that case with
$n = 2$. CARTA's algebra is stated for $n \ge 2$. Nothing in the reconstruction
problem requires exactly two fragments, and a user who suffers repeated device
failures can accumulate more. The generalisation runs one way: CARTA covers more
than TCF's minimal case, while TCF itself remains binary.

\subsection{Ingestion}
\label{sec:carta:ingestion}

Ingestion maps each fragment independently to a sequence of structured records:

\begin{equation}
  \mathcal{I}: F_k \mapsto R_k
  = \langle r_{k,1}, r_{k,2}, \dots, r_{k,M_k} \rangle
\end{equation}

Each record keeps the raw text slice, the raw timestamp string, the raw sender
token, a message category, and a provenance tuple:

\begin{equation}
  \mathrm{prov}(r_{k,j}) = \big\langle
    \texttt{source\_id}: F_k,\;
    \texttt{line\_range}: [\ell_{\text{start}}, \ell_{\text{end}}],\;
    \texttt{sha256}: h(F_k)
  \big\rangle
  \label{eq:prov}
\end{equation}

Three parts, three jobs. The source identifier says which fragment the record
came from. The line range makes the original text re-locatable, so a reader can
return to the bytes the record was derived from. The hash binds that claim to a
specific file state, which is what allows an auditor to detect that the fragment
has changed since the reconstruction ran.

Ingestion is non-lossy and non-destructive. It writes nothing back to the
fragment files, and it filters no content. It also performs no reconciliation
across fragments, and this restriction is deliberate rather than incidental: a
stage that both parses and reconciles would make it impossible to say afterwards
whether a discrepancy came from a parsing decision or a merging decision.
Ingestion parses. Reconciliation happens later, once, in one place.

\subsection{Normalization}
\label{sec:carta:normalization}

Normalization converts each fragment's records into a shared canonical schema,
with UTC timestamps, resolved sender identity, and a fixed message-type taxonomy:

\begin{equation}
  \mathcal{N}: R_k \mapsto (R'_k, Q_k)
\end{equation}

Each record meets one of three fates, and the third is the one that distinguishes
this stage from a conventional parser.

A record that conforms to the canonical grammar \textbf{unambiguously} becomes a
canonical record in $R'_k$, carrying a UTC epoch timestamp, a normalised sender
identity, and a typed payload. Unambiguous here means a date and time admitting a
single valid interpretation, valid UTF-8, and a sender or system token that
resolves.

A record that \textbf{fails to parse} moves to the quarantine set $Q_k$.

A record that \textbf{parses but admits more than one reading} also moves to
$Q_k$. The date string \texttt{03/04/24} is the standard case: it is well-formed
under both day-month-year and month-day-year conventions, and nothing inside the
record itself decides between them. Such a record is not normalised, and it is
not resolved by picking the more likely convention. Guessing here would convert
an uncertainty into an assertion that later stages, and eventually the reader,
would have no way to distinguish from an observation.

Every quarantined record carries a typed reason code $\rho$. The taxonomy
distinguishes at minimum \texttt{MALFORMED\_\allowbreak TIMESTAMP},
\texttt{UNSUPPORTED\_\allowbreak FORMAT}, \texttt{AMBIGUOUS\_\allowbreak DATE\_\allowbreak FORMAT},
\texttt{UNRESOLVED\_\allowbreak SENDER}, and \texttt{INVALID\_\allowbreak ENCODING}, which separates a
record that could not be read from one that could be read in several ways. The
distinction matters downstream, because the two call for different remedies: a
malformed record may need a parser fix, whereas an ambiguous one may be resolvable
by evidence outside the record, such as a locale established elsewhere in the same
fragment.

Normalization is required to carry every record's provenance through to $R'_k$ or
$Q_k$. This is a guarantee the stage owes, and not a property Synthesis may assume
on its behalf.

\subsection{Synthesis}
\label{sec:carta:synthesis}

Synthesis is the only stage that sees fragments together:

\begin{equation}
  \mathcal{S}: \big(\{R'_1, \dots, R'_n\}, \{Q_1, \dots, Q_n\}\big)
  \mapsto (\mathcal{U}, \Omega)
\end{equation}

It performs four operations on the union
$\mathcal{R}'_{\text{all}} = \bigcup_{k=1}^{n} R'_k$, in this order: it orders the
records chronologically, groups records that appear to describe the same message
into reconciliation clusters, synthesises one canonical record per cluster, and
detects observations that cannot be reconciled. Section~\ref{sec:algorithms}
gives the procedure.

Recall from Section~\ref{sec:tcf:conditions} that $\delta$ and $\epsilon$ may both
be strictly positive, so $F_1$ and $F_2$ can overlap. A message sent during that
overlap appears in both fragments and arrives at Synthesis twice, in two
separately parsed forms that need not agree byte for byte. Deduplication exists
because of this overlap. Without it, concatenation would report the overlap region
twice and misrepresent the history it claims to reconstruct.

A conflict is any set of two or more observations that cannot be reconciled
deterministically under the reconstruction rules. This is broader than a timestamp
collision, and it includes contradictory sender attribution and orderings that
cannot both hold.

\subsection{The reconstruction report}
\label{sec:carta:report}

$\Omega$ is a first-class output with four components:

\begin{equation}
  \Omega = \big\langle
    \Omega_{\text{provenance}},\;
    \Omega_{\text{quarantine}} = \textstyle\bigcup_k Q_k,\;
    \Omega_{\text{conflicts}},\;
    \Omega_{\text{stats}}
  \big\rangle
  \label{eq:omega}
\end{equation}

The \textbf{provenance map} records, for every record in $\mathcal{U}$, its origin
device, origin export, origin backup date, and confidence. The \textbf{quarantine
carryover} holds every record from every $Q_k$, retained with its reason code
rather than dropped at the stage boundary. The \textbf{conflict log} holds every
irreconcilable observation under the definition above. The \textbf{statistics}
report deduplication counts along with resolved and unresolved ambiguity counts.

Quarantine carryover is the component that does the most work. A reconstruction
that silently omits what it could not handle looks identical, to its reader, to
one that had nothing to omit. Carrying $\mathcal{Q}$ into $\Omega$ makes the two
cases distinguishable without requiring the reader to re-run anything.

\subsection{Design principles}
\label{sec:carta:principles}

\paragraph{Deterministic.}
Identical fragment inputs produce an identical $(\mathcal{U}, \Omega)$.

\paragraph{Auditable.}
Every record in $\mathcal{U}$ traces back to at least one source fragment through
its provenance.

\paragraph{Non-destructive.}
Original fragment files are read and never modified.

\subsection{Invariants}
\label{sec:carta:invariants}

Four invariants constrain any conforming implementation.

\paragraph{Invariant 1, Zero unaccounted loss.}
Every ingested record either contributes provenance to a record in the unified
archive or appears in the quarantine carryover:

\begin{equation}
  \forall r \in \bigcup_{k=1}^{n} R_k, \quad
  \big(\exists u \in \mathcal{U} : \mathrm{prov}(r) \subseteq \mathrm{prov}(u)\big)
  \;\lor\;
  \big(\exists q \in \Omega_{\text{quarantine}} : q = r\big)
  \label{eq:inv1}
\end{equation}

No third outcome is permitted. A record cannot leave the pipeline unaccounted for.

\paragraph{Invariant 2, Strict provenance traceability.}
Every record in the unified archive carries at least one provenance entry, and
every entry references a valid source fragment:

\begin{equation}
  \forall u \in \mathcal{U}, \quad
  |\mathrm{prov}(u)| \ge 1
  \;\land\;
  \forall p \in \mathrm{prov}(u),\; p \text{ references a valid } F_k
  \label{eq:inv2}
\end{equation}

This forbids synthesised records with no origin, which is what would let an
inference enter the archive indistinguishable from an observation.

\paragraph{Invariant 3, Determinism.}
For a fixed $\mathcal{F}$, two runs of the pipeline yield the same pair:

\begin{equation}
  \mathcal{S}(\mathcal{N}(\mathcal{I}(\mathcal{F}))) = (\mathcal{U}_1, \Omega_1)
  \;\land\;
  \mathcal{S}(\mathcal{N}(\mathcal{I}(\mathcal{F}))) = (\mathcal{U}_2, \Omega_2)
  \implies
  (\mathcal{U}_1 = \mathcal{U}_2) \land (\Omega_1 = \Omega_2)
  \label{eq:inv3}
\end{equation}

We state this as a requirement on conforming implementations. No implementation
exists to exhibit it, and Section~\ref{sec:algorithms:ordering} identifies the
tie-breaking rule that a conforming implementation needs in order to satisfy it.

\paragraph{Invariant 4, Idempotent reconstruction.}
Applying Synthesis to an already unified, canonical record set together with an
identical copy of itself must not introduce additional information loss or
unresolved conflicts. We state this as a specification-level target rather than a
proven equation. A type-safe formalisation, establishing that records in
$\mathcal{U}$ are schema-compatible with those in $\mathcal{R}'$, is deferred to
the implementation phase and is not claimed here.
